\chapter{Introducción}

Este informe documenta el desarrollo y la resolución del Ejercicio Práctico N° 4 del Trabajo Práctico N° 2.
El objetivo principal de esta actividad es la implementación de una interfaz de comunicación serie bidireccional
utilizando el periférico interno USART1 del microcontrolador STM32F103 (placa BluePill), comunicándolo con una PC
mediante un conversor físico USB-UART.

La comunicación serie asíncrona es uno de los métodos de comunicación más extendidos y sencillos en el ámbito de
los sistemas embebidos. En esta etapa, se diseña un sistema de control interactivo mediante una máquina de estados
que permite al usuario alternar entre modos de operación (como control de un parpadeo de LED y un modo de respuesta
asociado a comandos) desde una terminal en la PC. Para lograr esto, se configura el periférico USART1 en formato
9600 8N1, empleando interrupciones por hardware para transmisión (TC) y recepción (RXNE), garantizando un sistema
eficiente que no dependa de esperas activas prolongadas para la E/S serie.

Adicionalmente, se desarrolla una herramienta del lado de la PC en lenguaje C que actúa como monitor serie
interactivo por consola, facilitando las pruebas de transmisión, recepción y el envío de comandos hacia la placa.

